\section{Cross-Cutting Findings}\label{sec:crosscutting}

\subsection{Where we do well (high coverage and high agreement)}
\begin{itemize}[leftmargin=1.4em,itemsep=0.2em,topsep=0.2em]
\item Integer Sequences from Configurations in the Hausd (1997354) --- C=10, A=10
\item Mathematical Foundations of the GraphBLAS (1379592) --- C=9, A=10
\item Fortunato--Townsend ADI Poisson Solver --- C=9, A=10
\item Electronic/Optical Properties of 2D GaN (1484740) --- C=9, A=9
\item Pt-LTO Photocatalysis DFT (1981773) --- C=9, A=9
\item ScaWL k-WL Distributed (2587225) --- C=9, A=10
\item rVAE Parsimonious Representations (2439897) --- C=9, A=10
\item NukeLM Domain Language Models (1861801) --- C=8, A=10
\item Exactly-Solved Model of Light-Scattering Errors in (2441075) --- C=8, A=9
\item Clustering huge protein sequence sets in linear ti (1624105) --- C=8, A=9
\item Common Mode Voltage Reduction of Single-Phase Quas (1606674) --- C=8, A=9
\item NN-VMC A$\leq$4 Nuclei + 3-body (1864334) --- C=8, A=9
\item Markov State Models from Short Non-Equilibrium Tra (1565592-MSM-Hempel) --- C=9, A=8
\item NANOGrav 15-yr GWB, BLS Kepler, Mesh-GNN, ELM forecaster, fldgen, DMQMC+GPR, DRAS, NILC, CosmoPower (all C=8, A=8)
\item \textit{Wave~3 biology additions:} Price~2018 RB-TnSeq (C=10, A=9), Jiang~2017 ARG (C=9, A=9), Sivakumar~2023 \textit{S.\ aureus} (C=9, A=9), Thakur~2022 \textit{T.\ pyogenes} (C=9, A=9), Vassallo~2022 phage defence (C=8, A=9)
\item \textit{Wave~4 PDE additions:} Kochkov~2021 ML-CFD (C=8, A=9), Grossmann~2023 FEM-vs-PINNs (C=8, A=9)
\end{itemize}

\noindent\textbf{Patterns.} Papers at the top of the ranking share three
features: (a) the method is expressible in open-source Python/C++
(Markov-state models, topological responses, combinatorics, classical
graph algorithms, matched-filter cosmology, dark-matter direct detection
event rates, sequence alignment, learned PDE solvers); (b) the computational footprint fits on a
single workstation or single GPU node at paper scale; and (c) reference data
is public. In several cases (1997354 integer sequences, 1379592 GraphBLAS,
28589945 ARG identities) the replication recovered the paper's result
\emph{exactly}, including tables copied 1:1 or identity percentages matching
to 0.3\%. The Wave~4 results reinforce this pattern: the two REPLICATED
papers (Kochkov, Grossmann) both have public code and data, while the
three PARTIAL papers lack both.

\subsection{Where we do poorly (or honest negatives)}
\begin{itemize}[leftmargin=1.4em,itemsep=0.2em,topsep=0.2em]
\item Godunov-loss PDE (P1) --- C=4, A=8 (honest negative: MSE outperformed Godunov-hybrid loss for MLP)
\item Supervised extraction of near-complete genomes fro (2469515) --- C=3, A=5 (PATRIC pipeline proprietary)
\item Simple Coplanar Waveguide Resonator Mask Targeting (1983793) --- C=5, A=7
\item Generative AI-Driven Accelerated Discovery of Pass (PVMol-Gen-Fajar2026) --- C=7, A=5
\item A Portfolio Approach to Massively Parallel Bayesia (2571540) --- C=5, A=6
\item Cosmic Reionization On Computers (1275503) --- C=5, A=5
\item Divide and Conquer MP-NODE (3003857) --- C=5, A=4
\item \textit{Wave~3:} Centenarian bile acid (Sato 2021) --- C=4, A=6 (spot-check at 3--7\% read depth)
\item \textit{Wave~3:} Outer mucus niche (Li 2015) --- C=6, A=6 (Bray-Curtis not UniFrac; small effect sizes)
\item \textit{Wave~4:} Eivazi~2022 PINN-RANS --- C=6, A=5 (DNS/LES data unavailable; 0/14 quantitative claims reproduced)
\item \textit{Wave~4:} Zhang~2019 modal-space SPDE --- C=7, A=5 (no code; errors 20--35$\times$ higher; PINN training collapsed)
\end{itemize}

\noindent Four dominant failure modes explain the bottom quartile:
\begin{enumerate}
\item \textbf{Proprietary or gated pipelines.} BV-BRC/SEEDtk for supervised
genome binning (2469515), and CROC hydro codes for cosmic reionization
(1275503), required access we did not hold at time of evaluation. Access
requests are in flight.
\item \textbf{Compute scale.} A handful of papers require $10^4$--$10^5$
GPU-hours (6-band LSST latent SDE, full Kolmogorov MP-NODE, 8500-node
distribution grid). Our v1 attempts deliberately ran at reduced scale.
\item \textbf{Missing code + paywalled text.} 1868518 Graph-RL distribution
restoration had no public code and is behind a paywall; the replication is an
independent reimplementation with substantial inference gaps.
\item \textbf{PINN hyperparameter sensitivity.} Wave~4 demonstrates that PINN
papers without public code are systematically harder to reproduce: 3/3 PINN
papers without code achieved only PARTIAL status, with absolute errors
10--100$\times$ higher than claimed, despite correctly reproducing the
methodology and qualitative trends (see \S\ref{sec:pinn-repro}).
\end{enumerate}

\subsection{Common failure modes at finer grain}
\begin{itemize}
\item \textbf{Force-field / hyperparameter drift.} In MSM-Hempel (1565592),
our long-reference $t_2=2020$\,ps differs from the paper's $\sim$1400\,ps by
44\%, traceable to an unspecified force-field choice rather than a method bug.
\item \textbf{Tokenizer / representation drift.} In PVMol-Gen-Fajar2026, our
SELFIES-based generator outproduces the paper's SMILES pipeline 6.6$\times$ in
filtered yield, biasing downstream chemistry.
\item \textbf{Missing surrogate calculators.} Swapping xTB/DFT for RDKit
approximations in PVMol-Gen saved weeks but leaves electronic-property filters
unverified.
\item \textbf{Stochastic seeds.} Deep-RL experiments (1984484 DRAS, 1868518)
require many seeds; our v1 runs used only a few.
\end{itemize}

\subsection{Tool-substitution patterns}
Across the cohort, several recurring tool substitutions are documented:
\begin{itemize}[leftmargin=1.4em,itemsep=0.2em,topsep=0.2em]
\item \textbf{Bray-Curtis vs UniFrac}: microbiome papers (Li~2015) used
  Bray-Curtis as a proxy for weighted UniFrac, producing qualitatively
  similar ordination but quantitatively different effect sizes.
\item \textbf{abricate vs RGI}: AMR gene detection via abricate (CARD/ResFinder)
  vs the Resistance Gene Identifier yields equivalent gene calls at $\geq$95\%
  identity thresholds, with minor count differences.
\item \textbf{PLFam vs Roary}: BV-BRC's protein family clustering (PLFam)
  substitutes for Roary pan-genome analysis, producing consistent
  core/accessory partitions but different absolute gene counts (e.g., 3,412
  vs 4,360 for \textit{S.\ aureus}).
\item \textbf{FastANI vs pyani}: both ANI implementations agree within 0.1\%
  for conspecific comparisons; divergence appears only near the species
  boundary ($\sim$95\% ANI).
\end{itemize}

\subsection{PINN reproducibility}\label{sec:pinn-repro}

Three of the five Wave~4 papers (Eivazi~2022, Kopani\'cakov\'a~2023,
Zhang~2019) share a striking failure pattern: the \emph{methodology}
reproduces qualitatively (correct physics, correct algorithmic structure,
correct trend directions), but \emph{absolute numerical precision} does not.
Paper-claimed errors of 0.01--1\% become 1--50\% in our hands---a gap of
10--100$\times$.

The root causes are consistent across all three:
\begin{enumerate}[leftmargin=1.4em,itemsep=0.2em,topsep=0.2em]
\item \textbf{No public code.} All three papers lack published implementations.
  The two REPLICATED papers in this batch (Kochkov~2021, Grossmann~2023) both
  have public repositories, underscoring the importance of code availability.
\item \textbf{L-BFGS implementation sensitivity.} PINNs rely heavily on L-BFGS
  for the final refinement phase. Line-search strategy (cubic backtracking +
  strong Wolfe vs Armijo vs fixed step), history length, and convergence
  tolerances can shift final accuracy by 1--2 orders of magnitude. Standard
  library L-BFGS (PyTorch, scipy) does not match custom implementations.
\item \textbf{Undocumented hyperparameters.} Loss balancing coefficients,
  collocation point density, learning rate schedules, initialization seeds,
  and curriculum strategies are often omitted or incompletely specified. PINN
  training is known to be highly sensitive to these choices.
\item \textbf{Reference data availability.} Eivazi~2022 requires DNS/LES
  datasets from KTH that are not publicly available. Without the correct
  boundary data, the PINN cannot learn the correct field.
\end{enumerate}

\noindent\textbf{Contrast with the REPLICATED papers.} Kochkov~2021
(jax-cfd) provides open-source code, public training data on GCS, and
pre-trained model outputs---enabling independent verification. Grossmann~2023
provides a complete GitHub repository with all benchmark configurations.
Both achieved REPLICATED status despite being technically demanding. The
dichotomy is stark: \emph{code availability is the single strongest predictor
of replication success for PINN/ML-PDE papers.}

\noindent\textbf{Implication for the field.} The PINN reproducibility gap we
observe is consistent with findings by Krishnapriyan et al.\ (2021,
``Characterizing possible failure modes in physics-informed neural
networks'') and suggests that PINN papers should be held to a higher
standard of implementation disclosure than classical numerical methods,
precisely because their results are more sensitive to implementation
details.

\subsection{Spot-check vs replication boundaries}
The Wave~3 biology batch highlights the distinction between a \emph{spot-check}
and a \emph{replication}. Sato~2021 (centenarian bile acids, C=4/A=6) is
classified as SPOT-CHECK: at 3--7\% read depth, the analysis can confirm
directional trends but lacks statistical power for definitive claims. By
contrast, Jiang~2017 (ARG dissemination, C=9/A=9) is a full REPLICATION:
56/56 protein identities matched within $\pm$0.3\%. The BV-BRC API ceiling
also creates a natural boundary: the API works well when the paper's analysis
tools overlap with BV-BRC's capabilities (MLST, AMR gene calling, PLFam
pan-genome), but yields only PARTIAL when the paper relied on tools outside
BV-BRC's scope (Kleborate for serotyping, Roary for phylogenetics, DefenseFinder
for phage defence systems).
